\section{Reward Alignment with Stochastic Optimal Control}
\label{sec:appx:SOC}
In the reward alignment task for continuous-time generative models~\cite{uehara2024finetuningcontinuoustimediffusionmodels, domingo-enrich2025adjoint}, which our method builds upon, Uehara \textit{et al.}~\cite{uehara2024finetuningcontinuoustimediffusionmodels} introduce both an additional drift term $\psi$ (often referred to as a control vector field) and a modified initial distribution $\bar{p}_1$. 
Then the goal is to find $\psi$ and $\bar{p}_1$ such that the resulting final distribution at time $t=0$ matches the target distribution $p_0^*$ defined in Eq.~\ref{eq:problem definition and related work:target distribution at t=0}. Accordingly, the original reverse-time SDE used for generation (Eq.~\ref{eq:preliminaries:pretrained model SDE}) is replaced by a controlled SDE:
\begin{align}
    \label{eq:SOC:target SDE}
    \mathrm{d}\mathbf{x}_t &= \left(\mathbf{f}(\mathbf{x}_t, t)-\psi(\mathbf{x}_t, t) \right)\mathrm{d}t + g(t)\mathrm{d}\mathbf{W}, \qquad \mathbf{x}_1\sim \bar{p}_1.
\end{align}

This entropy-regularized stochastic optimal control framework adopts a pathwise optimization that integrates KL divergence penalties over trajectories and thus optimization formulation of Eq.~\ref{eq:problem definition and related work:optimization object at t=0} changes accordingly:
\begin{align}
    \label{eq:SOC:objective for controlled sde}
    \psi^*, p^*_1 = \argmax_{\psi, \bar{p}_1} \mathbb{E}_{\mathbb{P}^{\psi, \bar{p}_1}}\left[ r(\mathbf{x}_0) \right]-\alpha \mathcal{D}_{\text{KL}}\left[ \mathbb{P}^{\psi, \bar{p}_1} \|\mathbb{P}^{\text{data}} \right],
\end{align}
where $\mathbb{P}^{\psi, \bar{p}_1}$ is a measure over trajectories induced by the controlled SDE in Eq.~\ref{eq:SOC:target SDE} and $\mathbb{P}^{\text{data}}$ is a measure over trajectories induced by the pre-trained SDE in Eq.~\ref{eq:preliminaries:pretrained model SDE}.\\
KL-divergence term in Eq.~\ref{eq:SOC:objective for controlled sde} can be expressed as the sum of affine control cost $ \frac{\| \psi(\mathbf{x}_t, t) \|^2}{g^2(t)}$ and log Radon–Nikodym derivative at $t=1$, $i.e.$, $\log \frac{\bar{p}_1(\mathbf{x}_1)}{p_1(\mathbf{x}_1)}$:
\begin{align}
    \label{eq:SOC:objective for controlled sde girsanov}
    \psi^*, p^*_1 = \argmax_{\psi, \bar{p}_1} \mathbb{E}_{\mathbb{P}^{\psi, \bar{p}_1}}\left[ r(\mathbf{x}_0) \right]-\alpha\mathbb{E}_{\mathbb{P}^{\psi, \bar{p}_1}}\left[ \frac{1}{2} \int_{t=0}^1 \frac{\| \psi(\mathbf{x}_t, t) \|^2}{g^2(t)}\mathrm{d}t+\log \frac{\bar{p}_1(\mathbf{x}_1)}{p_1(\mathbf{x}_1)} \right],
\end{align}
which can be proved~\cite{uehara2024finetuningcontinuoustimediffusionmodels, domingo-enrich2025adjoint} using Girsanov theorem and martingale property of It$\hat{\text{o}}$ integral.\\
The optimal control $\psi^*$ and the optimal initial distribution $p_1^*$ can be derived by introducing the optimal value function, defined as:
\begin{align}
    \label{eq:SOC:optimal value funtion}
    V_t^*(\mathbf{x}_t)&=  \max_{\psi} \mathbb{E}_{\mathbb{P}^{\psi}}\left[ r(\mathbf{x}_0)- \frac{\alpha}{2}\int_{s=0}^t \frac{\| \psi(\mathbf{x}_s, s) \|^2}{g^2(s)}\mathrm{d}s \middle| \mathbf{x}_t \right].
    % &=\mathbb{E}_{\mathbb{P}^{*}}\left[r(\mathbf{x}_0)- \frac{\alpha}{2}\int_{s=0}^t \frac{\| \psi(\mathbf{x}_s, s) \|^2}{g^2(s)}\mathrm{d}s \middle| \mathbf{x}_t\right],
\end{align}
where the expectation is taken over trajectories induced by the controlled SDE in Eq.~\ref{eq:SOC:target SDE} with current $\mathbf{x}_t$ is given.

From the optimal value function at $t=1$, we can derive explicit formulation of the optimal initial distribution $p_1^*$ in terms of $V_1^*(\mathbf{x}_t)$ by plugging the definition of optimal value function at $t=1$ (Eq.~\ref{eq:SOC:optimal value funtion}) into Eq.~\ref{eq:SOC:objective for controlled sde girsanov}:
\begin{align}
    p_1^*=\argmax_{\bar{p}_1} \mathbb{E}_{\bar{p}_1}\left[ V_1^*(\mathbf{x}_1) \right]-\alpha \mathcal{D}_{\text{KL}}\left[ 
\bar{p}_1 \| p_1 \right].
\end{align}
Solving this yields the following closed-form expression for the optimal initial distribution (derivable via calculus of variations~\cite{kim2025inference}), similarly to Eq.~\ref{eq:problem definition and related work:target distribution at t=0}:
\begin{align}
    \label{eq:SOC:optimal initial distribution}
    p_1^*(\mathbf{x}_1)=\frac{1}{Z_1}p_1(\mathbf{x}_1)\exp\left( \frac{V_1^*(\mathbf{x}_1)}{\alpha} \right).
\end{align}

The optimal control $\psi^*$ can be obtained from the Hamilton–Jacobi–Bellman (HJB) equation and is expressed in terms of the gradient of the optimal value function:
\begin{align}
    \label{eq:SOC:optimal control}
    \psi^*(\mathbf{x}_t,t)=g^2(t)\nabla \frac{V_t^*(\mathbf{x}_t)}{\alpha}.
\end{align}
Moreover, the optimal value function itself admits an interpretable closed-form expression via the Feynman–Kac formula:
\begin{align}
    \label{eq:SOC:feynman-kac}
    V_t^*(\mathbf{x}_t)=\alpha \log \mathbb{E}_{\mathbb{P}^{\text{data}}}\left[ \exp \left( \frac{r(\mathbf{x}_0)}{\alpha} 
   \right) \middle| \mathbf{x}_t \right].
\end{align}\\
Importantly, Uehara \textit{et al.}~\cite{uehara2024finetuningcontinuoustimediffusionmodels} further proved that the marginal distribution $p_t^*(\mathbf{x}_t)$ induced by the controlled SDE (with optimal control $\psi^*$ and optimal initial distribution $p_1^*$) is given by:
\begin{align}
    \label{eq:SOC:optimal intermediate distribution}
    p_t^*(\mathbf{x}_t)=\frac{1}{Z_t}p_t(\mathbf{x}_t)\exp \left( \frac{V_t^*(\mathbf{x}_t)}{\alpha} \right),
\end{align}
where $p_t(\mathbf{x}_t)$ is the marginal distribution of the pretrained score-based generative model at time $t$.
Similarly, the optimal transition kernel under the controlled dynamics is:
\begin{align}
    \label{eq:SOC:optimal transition kernel}
    p_\theta^*(\mathbf{x}_{t-\Delta t}|\mathbf{x}_t)=\frac{\exp (V_{t-\Delta t}^*(\mathbf{x}_{t-\Delta t})/\alpha)}{\exp (V_t^*(\mathbf{x}_t)/\alpha)}p_\theta (\mathbf{x}_{t-\Delta t}|\mathbf{x}_t),
\end{align}
where $p_\theta (\mathbf{x}_{t-\Delta t}|\mathbf{x}_t)$ denotes the transition kernel of the pretrained model, $i.e.$, corresponding to the discretization of the reverse-time SDE defined in Eq.~\ref{eq:preliminaries:pretrained model SDE}.
For detail derivation, see Theorem 1 and Lemma 3 in Uehara \textit{et al.}~\cite{uehara2024finetuningcontinuoustimediffusionmodels}.
Notably, Eq.~\ref{eq:SOC:optimal intermediate distribution} implies that by following the controlled dynamics defined by Eq.~\ref{eq:SOC:target SDE}, initialized with the optimal distribution $p_1^*$ and guided by the optimal control $\psi^*$, the resulting distribution at time $t=0$ will match the target distribution $p_0^*$ defined in Eq.~\ref{eq:problem definition and related work:target distribution at t=0}.\\

Note that the optimal control, optimal initial distribution, and optimal transition kernel are all expressed in terms of the optimal value function. 
However, despite their interpretable forms, these expressions are not directly computable in practice due to the intractability of the posterior $p(\mathbf{x}_0|\mathbf{x}_t)$.
This motivates the use of approximation techniques, most notably Tweedie’s formula~\cite{efron2011tweedie}, which is widely adopted in the literature~\cite{chung2023dps, ye2024tfg, wu2023tds, kim2025DAS} to make such expressions tractable. 
Under this approximation, the posterior is approximated by a Dirac-delta distribution centered at the posterior mean denoted by $\mathbf{x}_{0|t}:=\mathbb{E}_{\mathbf{x}_0\sim p_{0|t}}[\mathbf{x}_0]$, representing the conditional expectation under $p_{0|t}:=p(\mathbf{x}_0|\mathbf{x}_t)$. 
Consequently, the optimal value function simplifies to:
\begin{align}
     V_t^*(\mathbf{x}_t)=\alpha \log \int \exp \left( \frac{r(\mathbf{x}_0)}{\alpha}\right)p(\mathbf{x}_0|\mathbf{x}_t)\mathrm{d}\mathbf{x}_0 \simeq \alpha \log \int \exp \left( \frac{r(\mathbf{x}_0)}{\alpha}\right)\delta(\mathbf{x}_0-\mathbf{x}_{0|t}) \mathrm{d}\mathbf{x}_0=r(\mathbf{x}_{0|t}),
\end{align}
where $\mathbf{x}_{0|t}$ is a deterministic function of $\mathbf{x}_t$. 
Using this approximation, we have following approximations for the optimal initial distribution $\tilde{p}_1^*$, the optimal control $\tilde{\psi}^*$, and the optimal transition kernel $\tilde{p_\theta}^*$, which are used throughout the paper:
\begin{align}
    \label{eq:SOC:approx optimal initial distribution}
    \tilde{p}_1^*(\mathbf{x}_1)&:=\frac{1}{Z_1}p_1(\mathbf{x}_1)\exp\left(\frac{r(\mathbf{x}_{0|1})}{\alpha} \right)\\
    \label{eq:SOC:approx optimal control}
    \tilde{\psi}^*(\mathbf{x}_t)&=g^2(t)\nabla \frac{r(\mathbf{x}_{0|t})}{\alpha}\\
    \label{eq:SOC:approx optimal transition kernel}
    \tilde{p}_{\theta}^*(\mathbf{x}_{t-\Delta t}|\mathbf{x}_{t})&=\frac{\exp (r(\mathbf{x}_{0|t-\Delta t})/\alpha)}{\exp (r(\mathbf{x}_{0|t})/\alpha)}p_\theta(\mathbf{x}_{t-\Delta t}|\mathbf{x}_{t}).
\end{align}\\
It is worth noting that sampling from the optimal initial distribution is essential to theoretically guarantee convergence to the target distribution Eq.~\ref{eq:problem definition and related work:target distribution at t=0}. Simply following the optimal control alone does not suffice and can in fact bias away from the target, a phenomenon known as the \textit{value function bias} problem~\cite{domingo-enrich2025adjoint}. 
To the best of our knowledge, this is the first work to explicitly address this problem in the context of inference-time reward-alignment with score-based generative models.